\section{Which servicing policy does the business sentence describe?}
\label{exec17:l2-policy}

A request to improve early servicing can leave the later action rule
unspecified. That omission is more consequential than a missing estimator
option. An initial review may change a later condition, and a subsequent
contact may respond to that condition. Comparing first-stage actions while
holding the later action fixed therefore differs from comparing complete
operating policies. The longitudinal formula developed earlier makes this
distinction precise; the following finite construction shows the numerical
consequence. It specializes the inherited sequential-intervention framework
of \citet{hernan2020} to an authored example, with no empirical claim about
how an actual servicing review affects hardship.

Let $A_0$ denote an initial action, $L_1$ a later adverse-state indicator,
and $A_1$ a follow-up action taken after observing $L_1$. Stipulate
$\Pr(L_1=1\mid\operatorname{do}(A_0=0))=1/4$ and
$\Pr(L_1=1\mid\operatorname{do}(A_0=1))=3/4$. The initial action has a
direct benefit but increases the chance of the costly later state in this
deliberately adverse model. Define net recovery, in hypothetical dollars, by
\begin{equation}
 Y(a,l,b)=100+20a-40l+30bl-10b.
 \label{exec17:l2-sequential-reward}
\end{equation}
The last term is the follow-up cost and is already included in the outcome.
There is no additional initial cost in this construction. The structural
model uses an independent uniform disturbance to generate $L_1$ at the
stated probabilities, and the outcome is then deterministic given the
actions and state. These declarations supply the causal law for the example;
observed historical conditional frequencies would not by themselves supply
the same law.

For a fixed initial action $a$, let $g^a$ take the follow-up action precisely
when $L_1=1$. In that state, follow-up adds USD20 net; when $L_1=0$ it
subtracts USD10. The adaptive rule is consequently optimal within this
specified two-action second-stage class when the state is fully observed
and there is no shared capacity constraint. Its value is
\begin{equation}
 V(g^a)=\sum_{l=0}^1Y(a,l,l)\Pr(L_1=l\mid\operatorname{do}(A_0=a)),
 \qquad V(g^0)=95,\quad V(g^1)=105.
 \label{exec17:l2-sequential-value}
\end{equation}
The complete-policy contrast is USD10. If follow-up is prohibited under
both initial actions, both values are USD90 and the first-action contrast
is zero. If the state is held fixed and the same follow-up action is used,
the direct contrast in the reward expression is USD20. Finally, comparing
$g^1$ with never taking either action yields USD15. All four calculations
are correct for their own interventions. A translator that reports one
without stating which subsequent policy is held fixed has not resolved the
business question. Subtracting the follow-up expense again would introduce
a separate accounting error.

The value comparison can be independently checked under a fully supported
generated logger. Let the logger randomize each action with probability
$1/2$, using an independent randomization at the second stage after observing
the state. Its eight possible histories have probabilities
$\tfrac14\Pr(L_1=l\mid\operatorname{do}(A_0=a))$. For the target $g^a$,
retain histories with $A_0=a$ and $A_1=L_1$ and multiply their logger
probabilities by four. Summing the weighted rewards reproduces 95 and 105.
This finite enumeration and the direct policy calculation are implemented
separately in the exact reference \citep{exec17_l2_policy_exact}. The agreement
checks this particular known law. It does not show that a real logger is
known, that hidden hardship is observed, or that weights fitted to dependent
customer histories have calibrated uncertainty.

The availability clock gives the rule an operational interpretation. If
$L_1$ is available before the second decision, $A_1=L_1$ is admissible.
If the note describing it arrives after that decision, the same expression
uses unavailable information. A fixed always-contact rule has values
USD87.50 and USD102.50 for initial actions zero and one; a fixed never-contact
rule has values USD90 and USD90. The best fixed rule is never-contact for initial action zero and
always-contact for initial action one under this stipulated law. These rules
illustrate a value difference caused by changing the permitted information,
not by switching estimators. A forecast available earlier is another
candidate input with its own joint law; it cannot be substituted for the
later realized state without a new calculation.

The formulation should accordingly distinguish three requests: evaluation
of a supplied policy, selection from a declared feasible policy class, and
an unrestricted optimal-policy claim. The example resolves only a tiny
known-law class. Adding hidden state, nonbinary actions, capacity constraints
or an unknown behavior policy creates additional research and implementation
requirements. The assistant can preserve these requirements in a formal
query and explain why the static effect is insufficient. It cannot convert
the finite optimum into a guarantee for general private-state servicing.

\section{A threshold change is an applicant-population decision}
\label{exec17:l2-threshold}

Credit approval introduces a different translation failure: observing a good
outcome among accepted applicants does not establish the value of expanding
acceptance. Let $D_z(S)$ be a deterministic approval rule at score $S$ under
policy $z$. Define $R(1)$ as net monetary contribution if a specified loan
is offered and $R(0)$ as the contribution under the specified alternative,
such as declining this product while retaining other services. On a common
applicant population, the policy value is
\begin{equation}
 V(z)=\mathbb E[D_z(S)R(1)+\{1-D_z(S)\}R(0)].
 \label{exec17:l2-credit-value}
\end{equation}
This definition needs the loan terms, horizon, default and recovery treatment,
capital charges and the decline alternative fixed. Repayment probability
alone does not supply $R$. Nor is the mean contribution among approved
applications equal to $V(z)$: its denominator changes with the threshold.

Suppose the new rule expands approval monotonically and changes no loan
terms for already approved applicants. Under a common potential-outcome law
without capacity spillovers, subtracting the policy values gives
\begin{equation}
 V(1)-V(0)=\mathbb E[\{D_1(S)-D_0(S)\}\{R(1)-R(0)\}].
 \label{exec17:l2-credit-switchers}
\end{equation}
Only the newly approved band contributes. If repricing or other changes also
affect incumbent approvals, this simplification fails and their contribution
must remain in the target. If actions interact through a capital or staffing
constraint, the individual response law also needs revision. The equation
is an accounting consequence of the specified potential outcomes, not an
identification result from historical approvals.

For an exact counterexample, let the old rule approve one fifth of applicants,
each with USD100 expected net contribution, and let the new rule admit an
additional one tenth. Set the decline contribution to zero. In one admissible
law the new band would contribute USD50 each if offered the loan; in another
it would contribute minus USD50. The historical data under the old rule can
be identical because this band received no loans. Its loan outcomes remain
unobserved. The old policy value is USD20 per applicant in either law, while
the expanded policy has values USD25 and USD15. The incremental values
are therefore USD5 and minus USD5. Exact enumeration verifies the result
\citep{exec17_l2_policy_exact}.

Meanwhile, mean contribution per approved applicant falls from USD100 to
USD$250/3$ in the favorable law, and to USD50 in the adverse law. Thus even
an observed reduction in the approved-group average does not alone show that
expansion reduces total applicant-population value. It can reflect adding a
positive but less profitable group. The proposal should carry the common
applicant denominator, the new approval mass and the net contribution
convention through the analysis. A question about whether loan quality falls
should retain portfolio value and conditional credit quality as
different questions until the owner chooses the intended objective.

Suppose only the scientific range
$-200\leq R(1)-R(0)\leq100$ is supplied for the new band. Its known mass
of one tenth then yields the incremental value range $[-20,10]$ dollars
per applicant. Both endpoints are attained by assigning the corresponding
constant contribution throughout the unobserved band, without changing the
old-policy observations. This is a sharp range within the stipulated class,
not a sampling interval or a calibrated statement about a bank's loss
distribution. A more restrictive response model, credible experimental
information or justified transport restriction can narrow it. A prediction
for the unobserved band merely names one extrapolation unless its scientific
basis is established.

A discontinuity at the old threshold can, under its design assumptions,
answer a local question about applicants near that threshold. It does not
automatically supply the response throughout a finite newly admitted band.
The assistant must distinguish moving the cutoff by a specified amount
from estimating a local effect at the existing cutoff, and retain the
support or transport argument required to connect them. Public threshold
mechanisms and generated response laws can test this distinction before a
bank considers a costly pilot. The relevant prototype success is correct
translation, appropriate bounds and justified abstention under absent
information, not confident completion of every numerical request.
\label{exec17:l2-policy-end}
